\chapter{Isu Implementasi Pengontrol Digital}
\label{bab:impl}
\enchap{Implementation Issues of Digital Controllers}

\begin{tujuan}
	Setelah mempelajari bab ini, Anda mampu menganalisis pengaruh tunda komputasi dan \textit{jitter}, menjelaskan efek kuantisasi koefisien dan aritmetika titik tetap, menyusun struktur program waktu nyata yang aman, dan menerapkan alur uji dari simulasi ke perangkat keras.
\end{tujuan}

\section{Pewaktuan}
\subsection{Tunda komputasi}
Dalam model diskrit standar, $u[k]$ dianggap berlaku tepat pada saat $y[k]$ dicacah. Pada kenyataannya aksi kontrol baru dikirim setelah waktu komputasi $T_c$. Ada tiga strategi: (1) kirim segera setelah dihitung sehingga $T_c \ll T$ dan dapat diabaikan atau diperlakukan sebagai tunda kecil; (2) kirim pada cacah berikutnya sehingga tunda tepat $z^{-1}$ dan dimasukkan ke model saat merancang; (3) pisahkan komputasi, yaitu hitung bagian yang hanya bergantung pada data lama sebelum cacah dan sisakan sedikit operasi setelah membaca $y$.

\begin{contoh}[Tunda satu cacah yang tidak dimodelkan]
	Kompensator \textit{lead} Bab~\ref{bab:rlocus} pada plant posisi memberi $M_p = 16\%$ dan $t_s = 2{,}1$\,s. Bila aksi kontrolnya ternyata dikirim satu cacah kemudian, lewatan naik menjadi $38\%$, waktu turun menjadi 3,8\,s, dan jari-jari kutub terbesar naik dari 0,82 menjadi 0,90. Tunda satu cacah menambah fasa $-\omega T$, dua kali tunda ZOH. Bila $u$ dikirim di cacah berikutnya, rancang ulang dengan $z^{-1}G(z)$.
\end{contoh}

\subsection{\textit{Jitter}}
\term{\textit{Jitter}}{jitter} adalah variasi waktu antar-cacah atau antar-keluaran, padahal model diskrit mengasumsikan $T$ konstan. Pada Arduino, sumbernya antara lain \texttt{Serial.print} yang panjang di dalam kalang, komunikasi I$^2$C dan \texttt{analogRead} berulang, cabang program dengan durasi berbeda, serta pemakaian \texttt{delay()} alih-alih pewaktuan \texttt{millis()}. Praktik yang baik adalah mengukur \texttt{DeltaTime} dan waktu eksekusi dengan \texttt{micros()} (Modul 0), menjaga waktu eksekusi terburuk di bawah $0{,}5T$, mencetak data lebih jarang atau dalam format biner ringkas, dan memakai \textit{timer interrupt} untuk $T$ kecil yang menuntut presisi.

\section{Numerik}
\subsection{Kuantisasi koefisien}
Koefisien kontroler disimpan dengan presisi terbatas. Bila kutub berada dekat $z = 1$ (karena $T$ kecil dibanding dinamika sistem), koefisien polinomial hampir saling meniadakan sehingga galat kecil menggeser akar jauh.

\begin{contoh}[Sensitivitas kutub]
	Filter orde dua dengan kutub $0{,}98e^{\pm j0{,}05}$ memiliki penyebut $z^2 - 1{,}957551z + 0{,}960400$. Bila koefisien dibulatkan ke 8 bit pecahan, kutub menjadi $0{,}98027e^{\pm j0{,}0599}$, yaitu frekuensi osilasi bergeser 20\%. Dengan 12 bit sudutnya 0,0508 rad, dan dengan 16 bit kembali 0,0500 rad.
\end{contoh}

Mitigasinya: realisasikan kontroler sebagai kaskade atau paralel seksi orde satu dan dua, bukan bentuk langsung orde tinggi; gunakan bentuk ruang keadaan atau operator delta $\delta = (z-1)/T$; dan jangan memilih $T$ jauh lebih kecil dari yang diperlukan.

\subsection{Titik tetap dan titik mengambang}
ATmega328P pada Arduino UNO tidak memiliki unit \textit{floating-point}; tipe \texttt{float} diemulasi perangkat lunak dan jauh lebih lambat daripada operasi bilangan bulat, sehingga waktu eksekusinya perlu diukur dengan \texttt{micros()}. Pada AVR, \texttt{double} sama dengan \texttt{float} 32-bit (sekitar 7 digit desimal). Aritmetika titik tetap (format Q) memakai bilangan bulat dengan skala $2^{-b}$; cepat, tetapi rentangnya terbatas sehingga \textit{overflow} harus diwaspadai dan aritmetika jenuh digunakan.

\begin{contoh}[Koefisien $a = e^{-T/\tau}$ dalam Q8]
	Dengan $\tau = 0{,}5$\,s dan 8 bit pecahan: untuk $T = 0{,}1$\,s, $a = 0{,}818731$ menjadi $0{,}820313$ dengan $\tau$ efektif $0{,}505$\,s; untuk $T = 0{,}01$\,s, $\tau$ efektif $0{,}507$\,s; untuk $T = 0{,}001$\,s, $a = 0{,}998002$ menjadi $0{,}996094$ dengan $\tau$ efektif hanya $0{,}255$\,s. Waktu cacah yang sangat kecil membuat model kehilangan separuh konstanta waktunya.
\end{contoh}

\section{Praktik implementasi}
\subsection{Struktur program}
Struktur program waktu nyata sebaiknya mengikuti urutan: ukur, hitung, batasi, kirim, lalu pekerjaan lain seperti memperbarui pengamat, \textit{anti-windup}, dan pencatatan data (Kode~\ref{lst:b14-loop}). Penulisan \texttt{last += T\_US} menjaga rata-rata periode tepat $T$ tanpa \textit{drift}. Aspek keselamatan meliputi batas $u$, deteksi sensor di luar rentang, dan perintah berhenti darurat. Semua keadaan (integrator, pengamat) diinisialisasi di \texttt{setup()} dan saat berganti mode.

\begin{lstlisting}[language=C++,caption={Kerangka kalang kontrol waktu nyata},label={lst:b14-loop}]
void loop() {
  baca_perintah_serial();          // non-blocking
  unsigned long now = micros();
  if (now - last >= T_US) {
    last += T_US;                  // tanpa drift
    float y = baca_sensor();       // 1. ukur
    float u = kontroler(r, y);     // 2. hitung
    u = batasi(u);                 // 3. saturasi
    run_motor(kompensasi(u));      // 4. kirim
    update_keadaan(u);             // 5. pengamat, anti-windup
    if (++n % 5 == 0) log_data();  // 6. catat jarang
  }
}
\end{lstlisting}

\subsection{Nonlinieritas}
Rancangan linier mengasumsikan beberapa nonlinieritas sudah ditangani: saturasi aktuator (batasi $u$ dan $\Delta u$, terapkan \textit{anti-windup} pada integrator dan pengamat dengan memakai $u$ setelah saturasi); \textit{deadzone} dan gesekan (kompensasi invers \textit{deadzone} seperti pada Modul 1); serta sensor (filter \textit{anti-alias}, kalibrasi, satuan yang konsisten, dan perilaku potensiometer yang melompat antara 100\% dan 0\%). Sebagian besar kegagalan implementasi berasal dari hal-hal ini, bukan dari teori kontrolnya.

\subsection{Alur uji}
Alur uji yang dianjurkan terdiri atas lima tahap: (1) simulasi model di MATLAB atau Python; (2) simulasi kode kontroler yang sama persis dengan yang akan diunggah; (3) uji pada kit tanpa daya motor sambil mengamati \textit{Serial Monitor} (SOP Modul 0: lepas transistor atau jangan sambungkan 220\,V); (4) uji pada kit dengan penguatan dan \textit{setpoint} kecil; (5) uji penuh termasuk gangguan, dengan pencatatan data. Tahap 2 menangkap kesalahan satuan, tanda, dan urutan operasi sebelum motor bergerak. Data kit dan simulasi sebaiknya dibandingkan pada grafik yang sama; bila berbeda jauh, periksa model (identifikasi ulang), tunda, saturasi, dan \textit{deadzone}.

\section{Peta metode}
Tabel~\ref{tab:b14-peta} merangkum metode yang dipelajari dalam kuliah ini. Semuanya bermuara pada persamaan beda yang dijalankan setiap $T$ detik di mikroprosesor.

\begin{table}[H]
	\centering
	\caption{Peta metode perancangan kontrol diskrit}
	\label{tab:b14-peta}
	\small
	\renewcommand{\arraystretch}{1.3}
	\begin{tabularx}{\textwidth}{@{}l l X l@{}}
		\toprule
		Pendekatan & Model & Metode & Bab \\
		\midrule
		Heuristik & FOPDT dari data & PID dengan ZN/Cohen--Coon ($\theta + T/2$) & \ref{bab:emulasi} \\
		Emulasi & $G(s)$ & Tustin/MPZ dari $C(s)$ & \ref{bab:emulasi} \\
		Langsung $z$ & $G(z)$ (ZOH) & \textit{root locus}, bidang $w$ & \ref{bab:rlocus}, \ref{bab:frek} \\
		Langsung $z$ & $G(z)$ & Ragazzini, \textit{deadbeat}, Dahlin & \ref{bab:frek} \\
		Ruang keadaan & $\Phi, \Gamma, C$ & penempatan kutub dan pengamat & \ref{bab:pp}, \ref{bab:obs} \\
		Optimal & $\Phi, \Gamma$ dan derau & LQR, Kalman, LQG, MPC & \ref{bab:lqr} \\
		\bottomrule
	\end{tabularx}
\end{table}

\begin{ringkasan}
	Tunda komputasi diatasi dengan mengirim aksi segera atau memodelkannya sebagai $z^{-1}$. Periode cacah dijaga tetap dan \textit{jitter} diukur. Koefisien kutub dekat $z = 1$ peka terhadap kuantisasi sehingga dipakai realisasi orde rendah. Implementasi diuji bertahap dengan menangani saturasi, \textit{deadzone}, dan sensor.\\[3pt]
	\en{Handle computational delay by early output or by modelling it as $z^{-1}$; keep the sampling period constant and measure jitter; use low-order realisations to limit coefficient sensitivity; test incrementally and handle saturation, deadzone, and sensor issues.}
\end{ringkasan}

\paragraph{Istilah kunci.} tunda komputasi \en{(computational delay)}, \textit{jitter}, kuantisasi koefisien \en{(coefficient quantisation)}, titik tetap \en{(fixed point)}, \textit{timer interrupt}, aritmetika jenuh \en{(saturating arithmetic)}.

\section*{Latihan}
\begin{enumerate}
	\item Ukur waktu eksekusi fungsi PID (Kode~\ref{lst:b5-pid}) di Arduino UNO dengan \texttt{float} dan dengan bilangan bulat Q15.
	\item Rancang ulang kompensator \textit{lead} Bab~\ref{bab:rlocus} dengan memperhitungkan tunda $z^{-1}$.
	\item Tunjukkan bahwa realisasi kaskade dua seksi orde satu kurang peka terhadap kuantisasi daripada bentuk langsung orde dua untuk kutub $0{,}98$ dan $0{,}97$.
	\item Susun daftar periksa (\textit{checklist}) keselamatan sebelum menjalankan kontroler baru pada kit.
\end{enumerate}
